% Part: sets-functions-relations
% Chapter: relations
% Section: equivalence-relations

\documentclass[../../../include/open-logic-section]{subfiles}

\begin{document}

\olfileid[fa-IR]{sfr}{rel}{eqv}

\olsection{روابط هم‌ارزی}

رابطهٔ همانی روی هر مجموعه، رابطه‌ای بازتابی، متقارن و تراگذری است. روابط~$R$
که هر سهٔ این ویژگی‌ها را دارند بسیار رایج‌اند.

\begin{defn}[رابطهٔ هم‌ارزی]
رابطهٔ $R \subseteq A^2$ که بازتابی، متقارن و تراگذری باشد
\emph{رابطهٔ هم‌ارزی} نامیده می‌شود. گفته می‌شود !!^{element}s $x$
و $y$ از~$A$ \emph{نسبت به $R$ هم‌ارزند} اگر~$Rxy$.
\end{defn}

روابط هم‌ارزی مفهوم \emph{ردهٔ هم‌ارزی} را پدید می‌آورند. یک رابطهٔ
هم‌ارزی دامنه را به رده‌های متفاوتِ یک افراز «تکه‌تکه می‌کند». درون هر
رده، همهٔ اشیا با یکدیگر در رابطه‌اند؛ و هیچ دو شیئی از رده‌های متفاوت
با یکدیگر در رابطه نیستند. گاهی سودمند است که دربارهٔ خودِ این رده‌ها
\emph{مستقیماً} سخن بگوییم. برای این منظور، تعریف زیر را معرفی می‌کنیم:

\begin{defn}\ollabel{def:equivalenceclass}
فرض کنید $R \subseteq A^2$ رابطه‌ای هم‌ارزی باشد. برای هر $x \in A$،
\emph{ردهٔ هم‌ارزیِ} $x$ در~$A$ مجموعهٔ $\equivrep{x}{R}
= \Setabs{y \in A}{Rxy}$ است. \emph{مجموعهٔ خارج‌قسمتِ} $A$ نسبت
به~$R$ عبارت است از
$\equivclass{A}{R} = \Setabs{\equivrep{x}{R}}{x \in A}$، یعنی مجموعهٔ
این رده‌های هم‌ارزی.
\end{defn}

نتیجهٔ بعدی با اثبات اینکه رده‌های هم‌ارزی به‌راستی بخش‌های یک افراز از~$A$
هستند، تعریف ردهٔ هم‌ارزی را موجه می‌سازد:

\begin{prop}
اگر $R \subseteq A^2$ رابطه‌ای هم‌ارزی باشد، آنگاه $Rxy$ اگر و تنها اگر
$\equivrep{x}{R} = \equivrep{y}{R}$.
\end{prop}

\begin{proof}
برای جهت چپ به راست، فرض کنید $Rxy$ و بگذارید $z \in
\equivrep{x}{R}$. بنا بر تعریف، آنگاه $Rxz$. چون $R$ رابطه‌ای هم‌ارزی
است، $Ryz$. (با شرح جزئیات: چون $Rxy$ برقرار است و~$R$ متقارن است، $Ryx$؛
و چون $Rxz$ برقرار است و~$R$ تراگذری است، داریم~$Ryz$.) پس
$z \in \equivrep{y}{R}$. با تعمیم
این استدلال، $\equivrep{x}{R} \subseteq \equivrep{y}{R}$. اما درست به
همین ترتیب، $\equivrep{y}{R} \subseteq \equivrep{x}{R}$. بنابراین بنا
بر اصل گسترش، $\equivrep{x}{R} =
\equivrep{y}{R}$.

برای جهت راست به چپ، فرض کنید $\equivrep{x}{R} =
\equivrep{y}{R}$. چون $R$ بازتابی است، $Ryy$، پس $y \in
\equivrep{y}{R}$. پس $y \in \equivrep{x}{R}$ نیز، بنا بر این فرض که
$\equivrep{x}{R} = \equivrep{y}{R}$. بنابراین $Rxy$.
\end{proof}

\begin{ex}
نمونه‌ای خوب از روابط هم‌ارزی از حساب پیمانه‌ای به دست می‌آید. برای هر
سه عدد صحیح مثبت $a$، $b$ و $n \in \PosInt$، می‌گوییم $a \equiv_n b$ اگر و تنها اگر
تقسیم $a$ بر~$n$ همان باقی‌مانده‌ای را بدهد که تقسیم $b$ بر~$n$ می‌دهد.
(به‌صورتی نمادین‌تر: $a \equiv_n b$ اگر و تنها اگر یک $k \in
\Int$ وجود داشته باشد که $a - b = kn$.) اکنون، $\equiv_n$ برای
هر~$n$ رابطه‌ای هم‌ارزی است. و دقیقاً $n$ ردهٔ هم‌ارزیِ متمایز
به‌وسیلهٔ~$\equiv_n$ پدید می‌آیند؛ یعنی $\equivclass{\Nat}{\equiv_n}$
دارای $n$ !!{element}s است. این رده‌ها عبارت‌اند از: مجموعهٔ اعدادی که
بر $n$ بدون باقی‌مانده بخش‌پذیرند، یعنی $\equivrep{0}{\equiv_n}$؛
مجموعهٔ اعدادی که تقسیمشان بر $n$ باقی‌ماندهٔ~$1$ می‌دهد، یعنی
$\equivrep{1}{\equiv_n}$؛ \ldots؛ و مجموعهٔ اعدادی که تقسیمشان بر~$n$
باقی‌ماندهٔ~$n-1$ می‌دهد، یعنی~$\equivrep{n-1}{\equiv_n}$.
\end{ex}

\begin{prob}
نشان دهید که $\equiv_n$ برای هر $n \in
\PosInt$ رابطه‌ای هم‌ارزی است و
$\equivclass{\Nat}{\equiv_n}$ دقیقاً $n$ عضو دارد.
\end{prob}

\end{document}
